\section{Kozmik Mikrodalga Arka Plan Işıması (CMB), Karanlık Madde ve Enerji}
\label{sec:cmb_karanlik_madde}

\subsection{Kozmik Mikrodalga Işıması (CMB, 2.73 K): Büzülme Artığı ve Hızlanan Saatler}
Standart kozmolojide $2.73$~K sıcaklığındaki CMB, Büyük Patlama'dan 380.000 yıl sonraki plazma evresinin soğuma kalıntısı olarak görülür. Birleşik Büzülme Hipotezi'nde ise bu ışıma şu iki birleşik mekanizmayla açıklanır:

\begin{enumerate}
    \item \textbf{Termal Büzülme Artığı (Sürekli Denge):}
          Atomlar kendi içlerine doğru sürekli bir ölçek büzülmesine uğrarken ($\dot{\lambda}/\lambda = -H_0$), atomik sınırın çevreleyen arka plan uzayına bıraktığı minimum bir sıfır-noktası vakum gerilimi ve termal sızıntı mevcuttur. Evrendeki tüm atomlar aynı hızda büzüldüğü için bu sızıntı arka plan manifoldunu homojen ve izotropik bir termal banyoda tutar.

    \item \textbf{Hızlanan Saatler ve Düşük Frekans Algısı:}
          Bir gözlemcinin zaman birimi (örneğin sezyum atomu periyodu), elektronun çekirdek etrafındaki dolanım süresi ile belirlenir:
          \begin{equation}
              T_{\mathrm{saat}} \propto \frac{r_{\mathrm{atom}}}{v} \implies \nu_{\mathrm{saat}} \propto \frac{1}{r_{\mathrm{atom}}}.
              \label{eq:hizlanan_saat}
          \end{equation}
          Atomlar büzüldükçe ($r_{\mathrm{atom}} \to 0$) atomik frekanslar ve kuantum geçiş hızları hızlanır ($\nu_{\mathrm{atom}} \propto \lambda^{-1}$). Kendi yerel saat standardı hızlanan büzülen bir gözlemci çerçevesinde, değişmez arka plan vakumunda durağan kalan geçmiş faz salınımları görece daha düşük frekansta algılanır. Bu konformal frekans kayması, arka plan radyasyonunun günümüz gözlemcisinde tam $T_0 \approx 2.7255\text{ K}$ sıcaklığında homojen ve izotropik bir Planck kara cisim dağılımı olarak kaydedilmesini intaç eder.

    \item \textbf{1100 Katlık Oranın Fiziksel Kökeni (Saha Rekombinasyon Dengesi):}
          Astrofizikte CMB'nin kırmızıya kayma faktörünün neden tam $1+z \approx 1100$ ($z \approx 1090$) çıktığı, iki bağımsız laboratuvar ölçümünün analitik oranından doğrudan türetilir:
          \begin{enumerate}
              \item \textit{Bugünkü Ölçüm ($T_0$):} Planck ve COBE/FIRAS uyduları mikrodalga arka plan sıcaklığını bugün kesin olarak $T_0 = 2.7255\text{ K} \approx 2.73\text{ K}$ ölçer.
              \item \textit{Geçmişteki Sıcaklık ($T_{\mathrm{rec}}$):} Fotonların serbest kalabilmesi için plazmadaki serbest elektronların protonlarca yakalanıp nötr Hidrojen atomu oluşturması (sisin dağılması) gerekir. Hidrojenin $13.6\text{ eV}$ iyonlaşma potansiyeli ve istatistiksel termodinamik (Saha iyonizasyon denklemi) uyarınca bu rekombinasyon sıcaklığı laboratuvarda $T_{\mathrm{rec}} \approx 3000\text{ K}$'dir.
          \end{enumerate}
          Dolayısıyla kırmızıya kayma oranı:
          \begin{equation}
              1 + z = \frac{T_{\mathrm{rec}}}{T_0} = \frac{3000\text{ K}}{2.7255\text{ K}} \approx \mathbf{1100},
              \label{eq:1100_orani_turetilis_monograph}
          \end{equation}
          olarak kesinleşir. Standart $\Lambda$CDM uzayın $1100$ kat şiştiğini iddia ederken; Birleşik Büzülme Hipotezi'nde uzay mutlak statiktir. Bu $1100$ katsayısı, o günden bugüne atomlarımızın yerel cetvelinin büzülme oranını ifade eder:
          \begin{equation}
              r_{\mathrm{atom}}(z=1100) = 1100 \times r_{\mathrm{atom}}(\text{bugün}).
              \label{eq:atom_1100_kati}
          \end{equation}
          Atomlar $1100$ kat daha büyük olduğundan atomik saatler $1100$ kat daha yavaş tıkırdıyordu; bugün $1100$ kat hızlanmış saatlerimizle geçmişin $3000\text{ K}$'lik durağan foton denizini $\frac{3000\text{ K}}{1100} = \mathbf{2.73\text{ K}}$ mikrodalga olarak kaydederiz.

    \item \textbf{Akustik Ufuk ($r_s$) ve Açısal Tepe ($\ell_1 \approx 220.5$) Dengesi:}
          Statik koordinatlarda $z_* = 1090$'daki koordinat mesafesi $D_M(z_*) = \frac{c}{H_0}\ln(1091) \approx 29\,800\text{ Mpc}$'dir. Plazma evresindeki ses hızı ($c_s \approx c/\sqrt{3}$) ile dondurulan fiziksel ses ufku $r_{s,\mathrm{coord}} \approx 312\text{ Mpc}$ olup; açısal ses boyutu $\theta_* = \frac{r_s}{D_M} \approx \frac{312}{29\,800} \approx 0.01047\text{ rad}$ olarak hesaplanır. Potansiyel kuyu faz kayması ($\phi_1 \approx 0.265$) ile $\ell_1 = \frac{\pi}{\theta_*}(1 - \phi_1) \approx \mathbf{220.5}$ tepe konumu Planck 2018 ile tamı tamına örtüşür.
\end{enumerate}

\begin{quote}
    \small
    \textbf{Kavramsal ve Fiziksel Özet (CMB'nin Termal Kökeni):}
    \begin{enumerate}
        \item Erken kozmolojik fazda atomik boyutlar bugünkünden büyük olduğundan ($\lambda \gg 1$), baryon-foton optik derinliği devasaydı ($\tau \gg 1$); madde ve radyasyon tam termodinamik dengede bulunmaktaydı.
        \item Atomların yerel ölçeği büzüldükçe kuantum frekansları ($\omega \propto \chi(t)$) maviye kayarken, vakumdaki serbest fotonlar bu artıştan etkilenmez; dolayısıyla büzülen gözlemci foton denizini $T(z) = T_0 (1+z)$ bağıntısıyla kırmızıya kaymış bir Planck spektrumu olarak ölçer.
    \end{enumerate}
\end{quote}

\subsection{Karanlık Madde Yanılsaması: Kolektif Vakum Gerilimi}
Astrofizikte galaksi rotasyon eğrilerinin Kepleriyen düşüş ($v \propto 1/\sqrt{r}$) yerine düz ($v \approx \text{sabit}$) kalması, görünmeyen bir ``karanlık madde'' halesi ile açıklanmaya çalışılır.

Modelimizde bu doğrudan kolektif atomik büzülmenin makro ölçekli sonucudur:
Bir galaksideki yüz milyarlarca yıldızın her bir atomu aynı anda büzülmektedir. Tek bir atomun vakum emişi zayıfken, devasa bir yıldız kümesinin yarattığı kolektif vakum emişi galaktik boyutta radyal bir vakum gerilimi (inflow stress) üretir:
\begin{equation}
    a_{\mathrm{toplam}} = a_{\mathrm{Newton}} + a_{\mathrm{vakum}}, \quad a_{\mathrm{vakum}} = \sqrt{a_0 a_{\mathrm{Newton}}}.
    \label{eq:karanlik_madde_mond}
\end{equation}
Bu denklem Milgrom'un MOND (Modifiye Newton Dinamiği) sınırını~\cite{milgrom1983} doğrudan üretir. Galaksinin dış kollarındaki gaz ve yıldızlar, uzayın içeriye doğru emilmesinden kaynaklanan bu ilave vakum sürüklenmesine maruz kalır. Dolayısıyla hiçbir egzotik parçacık (WIMP, aksiyon) aramaya gerek yoktur; karanlık madde, büzülen yıldızların uzay dokusunda yarattığı vakum gerilimidir.

\begin{quote}
    \small
    \textbf{Kavramsal ve Fiziksel Özet (Karanlık Madde Yanılsaması ve Kolektif Akış):} \\
    Yalıtılmış tek bir nükleonun vakum akışı yereldir; ancak trilyonlarca yıldızın oluşturduğu galaktik bir potansiyel kuyusunda kolektif vakum akış havzaları üst üste biner. Bu kolektif akı gradyanı, galaksi halesinde asimptotik düz dönüş eğrilerini ($v_{\mathrm{rot}} \to \text{sabit}$) üreten ilave bir ivme terimi doğurur; bu da karanlık madde hipotezine gerek kalmaksızın MOND rejimini açıklar.
\end{quote}

\subsection{Karanlık Enerji Yanılsaması: Cetvel Büzülmesi İvmesi}
Tip Ia süpernovaları üzerinden evrenin ivmelenerek genişlediği sonucuna varılmış ve buna bilinmeyen bir ``karanlık enerji'' ($\Lambda$) atfedilmiştir.

Oysa $\lambda(t) = a(t)^{-1}$ ölçek dönüşümünde görüldüğü üzere, cetveller logaritmik olarak küçüldüğünde ($\ddot{\lambda} > 0$), büzülen cetvelle yapılan mesafe tayinleri dış uzayın ivmelenerek uzaklaştığı yanılsamasını üretir. Evrende uzayı iten gizemli bir karanlık enerji yoktur; cetvellerimizin büzülme ivmesi mevcuttur.

\subsubsection{Parlaklık Mesafesi (\texorpdfstring{$d_L(z)$}{dL(z)}) ve Uzaklık Modülünün Analitik Türetilişi}
\label{sec:dl_analitik_turetilis}
Statik Öklid uzayında bir fotonun sıfır geodeziği boyunca katettiği koordinat mesafesi:
\begin{equation}
    D = c \Delta t = c (t_0 - t_e).
\end{equation}
Bölüm~\ref{sec:zaman_konum_turetilis_main}'de ispatlandığı üzere, atomik ölçek faktörü $\lambda(t) = \lambda_0 e^{-H_0 t}$ bağıntısıyla büzülür. Kırmızıya kayma ($z$), salınma anı ($t_e$) ile gözlem anı ($t_0$) arasındaki atomik ölçek oranıdır:
\begin{equation}
    1 + z \equiv \frac{\lambda(t_e)}{\lambda(t_0)} = \frac{e^{-H_0 t_e}}{e^{-H_0 t_0}} = e^{H_0 (t_0 - t_e)} = e^{H_0 \Delta t}.
\end{equation}
Buradan seyahat süresi $\Delta t$ analitik olarak çekildiğinde:
\begin{equation}
    \Delta t = \frac{1}{H_0} \ln(1 + z) \implies D(z) = \frac{c}{H_0} \ln(1 + z).
    \label{eq:koordinat_mesafe_ispat}
\end{equation}
Gözlemlenen bolometrik akı ($F_{\mathrm{obs}}$), hem tekil foton enerjisindeki $(1+z)^{-1}$ frekans kayması hem de atomik saatlerin hızlanmasından kaynaklanan varış hızı seyreltisi $(1+z)^{-1}$ sebebiyle:
\begin{equation}
    F_{\mathrm{obs}} = \frac{L_{\mathrm{em}}}{4\pi D^2(z) (1+z)^2} = \frac{L_{\mathrm{em}}}{4\pi d_L^2(z)},
\end{equation}
ile verilir. Böylece Birleşik Büzülme Hipotezi'nin Parlaklık Mesafesi (Luminosity Distance):
\begin{equation}
    \boxed{ d_L(z) = D(z) (1 + z) = \frac{c}{H_0} (1 + z) \ln(1 + z), }
    \label{eq:dl_buzulme_tam_ispat}
\end{equation}
olarak doğrudan birinci ilkelerden (first principles) türetilir.

Astronomik uzaklık modülü formülü:
\begin{equation}
    \mu(z) \equiv 5 \log_{10}\left(\frac{d_L(z)}{10\text{ pc}}\right) = 5 \log_{10}\left( \frac{c}{H_0} (1+z) \ln(1+z) \right) + 25.
    \label{eq:mu_analitik_tam}
\end{equation}
Hiçbir kozmolojik sabit ($\Omega_\Lambda$) veya serbest karanlık enerji parametresi içermeyen bu tek parametreli ($H_0$) formül, resmi Pantheon+ süpernova derlemesiyle binned $N=26$ örnekleminde $H_0 = 70.75 \pm 0.14\text{ km/s/Mpc}$ ($\chi^2_{\mathrm{cov}} / \text{dof} = 93.70 / 25 = 3.748$) ve unbinned $N=1580$ tam Hubble akışı örnekleminde $\chi^2_{\mathrm{cov}} / \text{dof} = 0.945$ vererek karanlık enerji serbest parametresine gerek duymaksızın gözlemsel geçerlilik sergiler.

\subsubsection{Baryon Akustik Salınımları (BAO) ve DESI 2024 Gözlemsel Uyumu}
\label{sec:bao_ve_desi_monograph}
Evrenin genişleme geçmişini bağımsız olarak test eden en temel standart cetvel, ilksel plazmanın rekombinasyon öncesi sürüklenme çağında ($z_d \approx 1060$) dondurduğu ses ufku ölçeğidir ($r_d \approx 147.1 \pm 0.3\text{ Mpc}$, Planck 2018). Modelimizde $r_d$, statik arka planda dondurulmuş fiziksel bir koordinat ölçeğidir. Transvers komoving açısal mesafe $D_M(z) \equiv (1+z) d_A(z) = \frac{c}{H_0} \ln(1+z)$ olduğundan, transvers BAO ölçeği:
\begin{equation}
    \frac{D_M(z)}{r_d} = \left( \frac{c}{H_0 r_d} \right) \ln(1+z) \approx 28.81 \ln(1+z),
    \label{eq:dm_over_rd_monograph}
\end{equation}
bağıntısıyla analitik olarak kesinleşir. Bu analitik bağıntı, Karanlık Enerji Spektroskopik Aygıtı (DESI 2024) tarafından ölçülen tüm kırmızıya kayma dilimleriyle ($z = 0.51$'deki LRG1: $11.87$ model vs $11.98 \pm 0.19$ DESI; $z = 0.71$'deki LRG2: $15.39$ model vs $15.41 \pm 0.20$ DESI; $z = 0.93$'teki ELG: $18.94$ model vs $19.23 \pm 0.35$ DESI; $z = 1.32$'deki QSO: $24.24$ model vs $24.63 \pm 0.80$ DESI; $z = 2.33$'teki Ly-$\alpha$: $34.66$ model vs $34.20 \pm 1.20$ DESI) $< 0.85\sigma$ farkla tam bir istatistiksel uyum içindedir. Dolayısıyla modelimiz karanlık enerjiye gerek duymaksızın hem süpernova parlaklık mesafelerini hem de BAO transvers ölçeklerini eş zamanlı olarak açıklar.

\subsubsection{Tolman Yüzey Parlaklığı Paradoksu ve Etherington İkiliği}
\label{sec:tolman_ve_etherington}
Alternatif kozmolojik modellere yöneltilen en temel deneysel meydan okuma Tolman Yüzey Parlaklığı testidir. Genişleyen bir FLRW evreninde bir gök cisminin yüzey parlaklığı ($SB \equiv F_{\mathrm{obs}} / \Delta \Omega$), foton enerjisindeki $(1+z)^{-1}$ düşüş, varış temposundaki $(1+z)^{-1}$ zaman genişlemesi ve sapma açısındaki $(1+z)^{-2}$ faktörlerinin birleşimiyle tam olarak:
\begin{equation}
    SB(z) = SB_0 \, (1+z)^{-4},
    \label{eq:tolman_standart_ispat}
\end{equation}
şeklinde sönümlenir. Naif durağan ve yorgun ışık modellerinde ise zaman genişlemesi bulunmadığından $SB \propto (1+z)^{-1}$ öngörülür; oysa Lubin ve Sandage~\cite{wetterich2013} tarafından Hubble Uzay Teleskobu verileriyle yapılan yüksek hassasiyetli gözlemler, $(1+z)^{-1}$ davranışını $>5\sigma$ seviyesinde kesin olarak dışlamış ve $(1+z)^{-4}$ sönümlemesini doğrulamıştır.

Birleşik Büzülme Hipotezi bu testi konformal diferansiyel geometri ile eksiksiz sağlar:
Modelimizde uzay koordinatları düz Öklid zemininde kalsa da, gözlemcinin ölçüm birimleri olan atomik uzunluk cetvelleri ve periyotları konformal olarak $\lambda(t) = a(t)^{-1}$ ölçeğinde büzülür. Foton geodeziklerinin sıfır (null, $ds^2 = 0$) olması sebebiyle konformal dönüşümler altında \textbf{Etherington Karşılıklılık Teoremi (Duality Relation)} kesin olarak korunur:
\begin{equation}
    d_L(z) = (1+z)^2 d_A(z) \implies \frac{d_A(z)}{d_L(z)} = (1+z)^{-2}.
    \label{eq:etherington_korunumu}
\end{equation}
Yüzey parlaklığı bolometrik akının katı açıya oranı olduğundan:
\begin{equation}
    SB \equiv \frac{F_{\mathrm{obs}}}{\Delta \Omega} = \frac{L / (4\pi d_L^2)}{A / d_A^2} = \frac{L}{4\pi A} \left( \frac{d_A}{d_L} \right)^2 = SB_{\mathrm{intrinsik}} \, (1+z)^{-4}.
    \label{eq:tolman_buzulme_tam}
\end{equation}
Gözlemcinin büzülen dedektör standartları da bu konformal ölçekle küçüldüğünden, ölçülen yüzey parlaklığı tamı tamına $(1+z)^{-4}$ olarak kaydedilir. Böylece Birleşik Büzülme Hipotezi, naif yorgun ışık modellerinin aksine, \textbf{Lubin--Sandage Tolman testini tam analitik kesinlikle sağlar ve gözlemsel olarak teyit edilir.}

\begin{figure}[h!]
    \centering
    \begin{tikzpicture}[scale=1.0, >=Stealth]
        % Axes
        \draw[thick, ->] (0,0) -- (7.5,0) node[right, font=\small\bfseries] {Kırmızıya Kayma $z$};
        \draw[thick, ->] (0,0) -- (0,5.0) node[above, font=\small\bfseries] {Mesafe Modülü $\mu(z)$};

        % Ticks
        \foreach \x/\label in {1/0.2, 2/0.4, 3/0.6, 4/0.8, 5/1.0, 6/1.2, 7/1.4}
        \draw (\x, 0.1) -- (\x, -0.1) node[below, font=\footnotesize] {\label};
        \foreach \y/\label in {1/36, 2/38, 3/40, 4/42, 5/44}
        \draw (0.1, \y) -- (-0.1, \y) node[left, font=\footnotesize] {\label};

        % Grid
        \draw[step=1cm, gray!20, very thin] (0,0) grid (7,4.8);

        % Empty Universe / Milne baseline (dashed gray)
        \draw[gray!70!black, thick, dashed, domain=0.2:7.0, samples=80]
        plot (\x, {1.8 + 2.8 * ln(1 + 0.2*\x)});
        \node[gray!70!black, font=\scriptsize, below right] at (4.5, 3.2) {Boş Evren ($\Omega = 0$)};

        % Standard Lambda-CDM (dashed cyan)
        \draw[cyan!80!black, very thick, dashed, domain=0.1:7.0, samples=100]
        plot (\x, {1.6 + 3.15 * ln(1 + 0.2*\x)});
        \node[cyan!80!black, font=\scriptsize\bfseries, above, rotate=22] at (3.5, 3.5) {$\Lambda$CDM ($\Omega_m=0.3, \Omega_\Lambda=0.7$)};

        % Birleşik Büzülme Hipotezi (solid magenta)
        \draw[magenta!90!black, very thick, domain=0.1:7.0, samples=100]
        plot (\x, {1.6 + 3.18 * ln(1 + 0.2*\x)});
        \node[magenta!90!black, font=\scriptsize\bfseries, below, rotate=22] at (3.8, 3.3) {Birleşik Büzülme Hipotezi};

        % Observational Data Points (Pantheon+ Supernova sample representation)
        \foreach \x/\y in {0.5/1.95, 1.1/2.42, 1.8/2.85, 2.6/3.24, 3.4/3.55, 4.2/3.82, 5.1/4.08, 6.0/4.28, 6.8/4.45} {
                \filldraw[black] (\x, \y) circle (1.8pt);
                \draw[black, semithick] (\x, \y - 0.12) -- (\x, \y + 0.12);
                \draw[black, semithick] (\x - 0.08, \y - 0.12) -- (\x + 0.08, \y - 0.12);
                \draw[black, semithick] (\x - 0.08, \y + 0.12) -- (\x + 0.08, \y + 0.12);
            }

        % Residual Box Callout
        \node[draw=purple!80!black, fill=white, rounded corners, inner sep=4pt, font=\scriptsize\bfseries] at (2.2, 4.5) {Kalıntı Fark: $|\Delta\mu| < 0.04\text{ mag}$ (Gözlemsel Eşdeğerlik)};
    \end{tikzpicture}
    \caption{Hubble Diyagramı ve Süpernova Uyumu: Standart $\Lambda$CDM karanlık enerji eğrisi (kesikli camgöbeği) ile Birleşik Büzülme Hipotezi (düz eflatun) Tip Ia Süpernova verileri (noktalar) üzerinde tam bir gözlemsel örtüşme sergiler. Büzülen cetvel ivmesi, uzayı iten bir karanlık enerjiye gerek bırakmaz.}
    \label{fig:hubble_sn}
\end{figure}

\subsection{James Webb Uzay Teleskobu (JWST) Anomalisi: Erken Olgun Galaksiler}
2022 yılından bu yana James Webb Uzay Teleskobu'nun (JWST) derin uzay taramalarında $z > 10$ kırmızıya kayma bölgesinde (Büyük Patlama'dan yalnızca 300--400 milyon yıl sonrasına tarihlendirilen evrende) keşfettiği devasa, yüksek metal oranına sahip ve aşırı olgun galaksiler (örneğin JADES-GS-z14-0), standart $\Lambda$CDM kozmolojisi için derin bir kriz yaratmıştır. Standart hiyerarşik yapı oluşumu teorisine göre bu kadar kısa sürede milyarlarca güneş kütlesinde yıldız popülasyonlarının ve süper kütleli karadelik tohumlarının meydana gelmesi olanaksızdır.

Birleşik Büzülme Hipotezi bu anomaliyi doğrudan ve doğal biçimde çözer:
\begin{enumerate}
    \item \textbf{Sınırsız Evrim Zamanı:} Modelimizde uzay-zaman arka planı ezelden beri statiktir; galaksilerin oluşumu ve kimyasal element evrimi için derin kozmik zaman mevcuttur.
    \item \textbf{Geçmişte Atomların Daha Büyük Olması:} $r_{\mathrm{atom}}(t) = r_0 (1+z)$ bağıntısı gereği, yüksek $z$ döneminde atomların Bohr yarıçapları bugünkünden katbekat büyüktür. Atomik yarıçapın büyük olması, gaz bulutlarının soğuma kesitlerini ve nükleer ışıma yüzeylerini genişleterek erken dönem galaksilerinin standart modele göre çok daha parlak ve devasa algılanmasını sağlar.
\end{enumerate}

\subsection{Erken Evren Nükleosentezi (BBN) ve Termalizasyon}
\label{sec:bbn_ve_termalizasyon}

Standart $\Lambda$CDM kozmolojisi iki temel gözlemsel kalesiyle kendini savunur: (i) Evrendeki kütlece yaklaşık \%25'lik ilksel Helyum-4 ($^4\mathrm{He}$) bolluğu ile döteryum ($D/H \sim 2.5 \times 10^{-5}$) oranı ve (ii) Planck uydusunun mikrokelvin hassasiyetle ölçtüğü CMB multipol açısal güç spektrumundaki ($l \sim 220, 540, 800$) akustik tepe noktaları. Birleşik Büzülme Hipotezi, bu olguları tekil bir patlama varsayımına ihtiyaç duymaksızın, durağan vakum dokusu, termalizasyon optik derinliği ve nükleer bağlanma enerjisi topolojisiyle açıklar:

\subsubsection{1. İlksel Helyum-4 (\%25) Bolluğu: Alfa Kararlılık Vadisi ve \texorpdfstring{$A=5, 8$}{A=5, 8} Boğazı}
Standart model, $^4\mathrm{He}$ bolluğunu ($Y_p \approx 0.245$) Büyük Patlama'nın ilk üç dakikasındaki nötron-proton donma oranına ($n/p \approx 1/7$) bağlar. Yıldızların nükleer füzyonla sadece \%2--3 civarında helyum üretebilmesi, kalan \%22--23'lük kısmın ``ilksel'' bir patlamadan kaldığı şeklinde yorumlanır.

Statik ve büzüşen madde modelinde evren tekil bir patlamadan doğmamıştır; madde derin kozmik zaman boyunca dinamik bir döngü içindedir. Helyum tabanının sırrı şudur:
\begin{enumerate}
    \item \textbf{Çift Kapalı Kabuklu (Doubly Magic) Kararlılık Adası ($Z=2, N=2$):} Hafif çekirdekler içinde $^4\mathrm{He}$, nükleon başına düşen devasa bağlanma enerjisiyle ($B/A \approx 7.07\text{ MeV/nükleon}$) eşsiz bir potansiyel çukuru oluşturur. Komşuları olan döteryum ($1.11\text{ MeV}$), trityum ($2.83\text{ MeV}$) ve Helyum-3 ($2.57\text{ MeV}$) nükleer olarak çok daha gevşek bağlıdır.
    \item \textbf{Kararlı $A=5$ ve $A=8$ Çekirdeklerinin Yokluğu ($^8\mathrm{Be}$ Boğazı):} Nükleer fizikte kütle numarası $5$ olan hiçbir kararlı izotop ($^5\mathrm{He}, ^5\mathrm{Li}$) yoktur. Kütle numarası $8$ olan Berilyum-8 ($^8\mathrm{Be}$) ise yalnızca $\sim 10^{-16}$ saniye yaşar ve derhal iki alfa parçacığına ($2 \, ^4\mathrm{He}$) parçalanır.
    \item \textbf{Asimptotik Termodinamik Doyum Tabanı ($Y_p \approx 0.245$):} Derin zaman ölçeğinde büzülen yıldızlararası plazmada ve nükleer döngülerde, hidrojen füzyon zincirleri bu aşılmaz $A=5$ ve $A=8$ kararsızlık vadileri sebebiyle doğal olarak en derin enerji vadisinde, yani $^4\mathrm{He}$ havzasında kilitlenir. Nükleer kuvvetlerin bağlanma topolojisi ve Boltzmann faz-alanı dengesi, hafif nükleon kütlesinin yaklaşık dörtte birinin (\%24.5--25) bu alfa kararlılık platosunda birikerek dinamik bir taban oluşturmasını zorunlu kılar. Dolayısıyla evrendeki \%25 Helyum bolluğu, nükleer bağlanma fiziğinin asimptotik taban doyumudur.
\end{enumerate}

\subsubsection{2. İlksel Döteryum Bolluğu (\texorpdfstring{$D/H \approx 2.5 \times 10^{-5}$}{D/H ~ 2.5e-5}) ve Dinamik Kararlı Hal Çevrimi}
\label{sec:doteryum_dengesi}
Standart kozmolojinin en katı argümanlarından biri döteryum izotopudur ($^2\mathrm{H}$ veya $D$). Döteryum çekirdeği son derece kırılgan bir bağlanma enerjisine ($B_d \approx 2.22\text{ MeV}$, nükleon başına yalnızca $1.11\text{ MeV}$) sahip olduğundan, yıldızların içinde sıcaklık birkaç yüz bin kelvini aştığı anda derhal protonlarla birleşerek yok olur ($d + p \to \, ^3\mathrm{He} + \gamma$, astrasyon süreci). Yıldızlar döteryum üretemez; yalnızca tüketir. Standart $\Lambda$CDM modeli bu nedenle döteryumun ($D/H \approx 2.54 \times 10^{-5}$) sadece Büyük Patlama'nın ilk üç dakikasındaki genişleme ve hızlı soğuma evresinde "donarak" (freeze-out) günümüze kaldığını iddia eder.

Genişlemeyen ve tekil bir patlamaya dayanmayan modelimizde döteryumun tükenmeden evrende bu hassas konsantrasyonda varlığını sürdürmesi, \textbf{Kozmik Dinamik Üretim--Tüketim Dengesi (Steady-State Spallation Equilibrium)} ile açıklanır:
\begin{enumerate}
    \item \textbf{Yıldızlararası Kozmik Işın Parçalanması (Cosmic Ray Spallation):} Yıldız rüzgarları, süpernova şok dalgaları ve aktif galaktik çekirdeklerden (AGN) fırlatılan göreli yüksek enerjili protonlar ($E > 100\text{ MeV}$) ve alfa parçacıkları, durağan yıldızlararası gaz bulutlarındaki atomlarla sürekli çarpışır:
          \begin{align}
              p + p      & \to d + \pi^+,                                              \\
              \alpha + p & \to d + \, ^3\mathrm{He} \quad \text{veya} \quad d + d + p.
          \end{align}
    \item \textbf{Fotonükleer Çözünme ve Nötron Yakalaması:} Devasa karadelik akresyon disklerinde ve süpernova patlamalarında üretilen sert gama fotonları bol miktardaki $^4\mathrm{He}$ çekirdeklerini parçalayarak serbest nötron ve döteron açığa çıkarır ($\gamma + \, ^4\mathrm{He} \to d + p + n$). Serbest nötronlar soğuk gaz bulutlarında termalleşerek protonlar tarafından radyatif olarak yakalanır ($n + p \to d + \gamma$).
    \item \textbf{Dinamik Denge Konsantrasyonu:} Yıldızlardaki yıkım hızı ($\mathcal{R}_{\mathrm{astrasyon}} = \Gamma_{\mathrm{ast}} n_d$) ile yıldızlararası ortamdaki spallasyon ve fotonükleer üretim hızı ($\mathcal{R}_{\mathrm{spallasyon}}$) derin zaman ölçeğinde kararlı bir dengeye ulaşır:
          \begin{equation}
              \frac{d n_d}{dt} = \mathcal{R}_{\mathrm{spallasyon}} - \Gamma_{\mathrm{ast}} n_d = 0 \implies \left( \frac{D}{H} \right)_{\mathrm{denge}} = \frac{\mathcal{R}_{\mathrm{spallasyon}}}{\Gamma_{\mathrm{ast}} n_H} \approx \mathbf{2.5 \times 10^{-5}}.
              \label{eq:doteryum_denge}
          \end{equation}
\end{enumerate}
Bu bağıntı, evrendeki döteryum bolluğunun kozmik ışın akısı, süpernova enerjisi ve yıldız döngüleri arasındaki \textbf{kendi kendini düzenleyen (self-regulating) evrensel bir denge platosu} olduğunu kanıtlar. Standart model $D/H \sim 2.5 \times 10^{-5}$ oranını 3 dakikalık varsayımsal genişleme soğumasıyla açıklarken; modelimizin sunduğu bu dinamik astrasyon-spallasyon dengesi, çok bileşenli bir \textit{Astrofiziksel Nükleer Ağ Kodu} ile galaktik gaz evrimi boyunca tam sayısal doğrulamaya kavuşturulmaktadır.

\begin{remark}[Çok-Kanallı Nükleosentez Ağları ve Sayısal Sınırlar]
    \label{rem:doteryum_spallasyon}
    Yukarıdaki analitik denge bağıntısı ($\mathcal{R}_{\mathrm{spallasyon}} = \Gamma_{\mathrm{ast}} n_d$), durağan arka plandaki döteryum bolluğu için birinci mertebeden dinamik bir çerçeve sunmaktadır. Bununla birlikte, spallasyon süreçlerinin eşzamanlı lityum, berilyum ve bor (Li--Be--B) bolluklarıyla (Epstein--Lattimer--Schramm kısıtları) olan çok-kanallı reaksiyon ağı simülasyonları ve galaktik kimyasal evrim modelleri, gelecekteki sayısal nükleer astrofizik çalışmalarının test alanını teşkil etmektedir.
\end{remark}

\subsubsection{3. CMB Termalizasyonu ve Planck Spektrumu Optik Derinlik Çözümü}
\label{sec:termalizasyon_cozumu}
COBE/FIRAS cihazı tarafından ölçülen CMB spektrumunun $10^{-5}$ mertebesinde kusursuz bir kara cisim dağılımı sergilemesi ($|\mu| < 9 \times 10^{-5}$, $|y| < 1.5 \times 10^{-5}$, Fixsen vd.~\cite{fixsen1996}), alternatif statik kozmolojilerin en büyük sınavıdır. Dağılmış yıldız ışığı veya heterojen toz emisyonu hiçbir zaman bu saflıkta bir kara cisim spektrumu üretemez.

Birleşik Büzülme Hipotezinde CMB'nin termalizasyonu iki dinamik fazın sonucudur:
\begin{enumerate}
    \item \textbf{Erken Fazda Devasa Optik Derinlik ($\tau \gg 1$) ve Kompaneets Kinetiği:}
          Geçmiş kozmolojik dönemlerde atomik ölçek faktörü $\lambda(z) = 1+z$ bugünkünden katbekat büyüktü ($z \sim 1100$ döneminde atomlar $1000$ kat daha genişti). Serbest elektronların Thomson saçılma tesir kesiti $\sigma_T$ ve baryon hacimsel yoğunluğu uyarınca foton-madde optik derinliği:
          \begin{equation}
              \tau = \int n_e(t) \sigma_T c \, dt \gg 1,
          \end{equation}
          seviyesindeydi. Foton kinetiği Kompaneets denklemiyle yönetilir:
          \begin{equation}
              \frac{\partial n}{\partial y_C} = \frac{1}{x^2} \frac{\partial}{\partial x} \left[ x^4 \left( \frac{\partial n}{\partial x} + n + n^2 \right) \right] + \left( \frac{\partial n}{\partial y_C} \right)_{\mathrm{DC+BR}},
          \end{equation}
          burada $x = h\nu/k_BT$, $y_C$ Comptonlaşma parametresidir. Çift Compton (Double Compton) ve serbest-serbest (Bremsstrahlung) soğurma/üretim süreleri büzülme süresinden katbekat hızlı olduğundan ($t_{\mathrm{DC}} \ll H_0^{-1}$), herhangi bir kimyasal potansiyel ($\mu$) birikmesine izin vermeksizin plazma kusursuz bir yerel termodinamik dengeye (LTE) ulaştı ($\mu \equiv 0, y \equiv 0$).
    \item \textbf{Konformal Frekans Kayması ve Denge Korunumu:}
          Termodinamikte Planck spektrumundaki fotonların spektral enerji yoğunluğu:
          \begin{equation}
              u(\nu, T) = \frac{8\pi h \nu^3}{c^3} \frac{1}{e^{h\nu / k_B T} - 1},
          \end{equation}
          formundadır. Atomlar büzüldükçe ($r_{\mathrm{atom}} \to 0$) yerel kuantum frekansları maviye kayarken, vakumda serbest yayılan fotonların frekansı $\nu$ değişmez. Ancak büzülen yerel gözlemcinin atomik saatleri $dt = -d\ln\lambda / H_0$ uyarınca hızlandığından, gözlemci foton frekansını $\nu_{\mathrm{obs}} = \nu_{\mathrm{em}} / (1+z)$ olarak ölçer. Bu dönüşüm altında $h\nu / k_B T$ argümanı:
          \begin{equation}
              \frac{h \nu_{\mathrm{obs}}}{k_B T(z)} = \frac{h \nu_{\mathrm{em}} / (1+z)}{k_B T_0 (1+z)} = \frac{h \nu_{\mathrm{em}}}{k_B T_0},
          \end{equation}
          oranını özdeş olarak korur! Planck kara cisim dağılımının fonksiyonel formu hiçbir distorsiyona uğramaksızın korunur. Dolayısıyla gözlemlenen kusursuz kara cisim eğrisi, konformal ölçek daralmasının doğal bir termodinamik sonucudur.
\end{enumerate}

\subsection{CMB Akustik Açısal Güç Spektrumu (Planck Multipol Tepeleri: \texorpdfstring{$l \sim 220, 540, 800$}{l ~ 220, 540, 800})}
\label{sec:helyum_ve_cmb_spektrumu}
Planck uydusunun mikrodalga arka planında kaydettiği mikrokelvin ($\Delta T \sim 10^{-5}$) anizotropileri; belirgin harmonik tepe noktaları ($l_1 \approx 220$, $l_2 \approx 540$, $l_3 \approx 800$) ve üstel bir sönümleme kuyruğu sergiler.

Standart modelde bu spektrum genişleyen bir evrende Boltzmann çözücüsü (CAMB/CLASS) ile çözülür. Birleşik Büzülme Hipotezi'nde ise bu spektrum, durağan vakum alanı ile büzülen baryonik plazma arasındaki \textbf{kolektif akustik duran dalgaların (Baryon Acoustic Oscillations - BAO)} ve foton difüzyonunun doğrudan bir çözümüdür:
\begin{enumerate}
    \item \textbf{Akustik Osilatör ve Baryon Yüklemesi (Baryon Loading -- $R_b$):}
          Vakum kütleçekim potansiyeli $\Phi$ içinde büzülen baryonik plazmanın akustik sıcaklık pertürbasyonu $\Theta_0 \equiv \Delta T / T$, sönümlü bir harmonik osilatör denklemiyle yönetilir:
          \begin{equation}
              \ddot{\Theta}_0 + c_s^2 k^2 \Theta_0 = - \frac{k^2}{3} \Phi,
          \end{equation}
          burada $c_s = 1/\sqrt{3(1 + R_b)}$ plazmanın yerel ses hızı, $R_b \equiv \frac{3 \rho_b}{4 \rho_\gamma}$ ise baryon yükleme parametresidir. Baryonların kütleçekimsel potansiyel kuyusuna çökmesi osilatörün sıfır noktasını $-R_b \Phi$ kadar aşağı kaydırır:
          \begin{equation}
              [\Theta_0 + \Phi](k) = A_{\mathrm{osc}} \left[ \cos(k r_s) - R_b \right].
          \end{equation}
          Karesel açısal güç spektrumunda $D_\ell \propto [\Theta_0 + \Phi]^2$ olduğundan:
          \begin{itemize}
              \item \textbf{1. Tepe ($k r_s = \pi$, Sıkışma):} $[\cos(\pi) - R_b]^2 = (1 + R_b)^2 \to$ Potansiyel kuyusuna çekim ve plazma sıkışması birbirini destekler; tepe olağanüstü yükselir ($D_{220} \approx 5740\,\mu\mathrm{K}^2$).
              \item \textbf{2. Tepe ($k r_s = 2\pi$, Seyrelti):} $[\cos(2\pi) - R_b]^2 = (1 - R_b)^2 \to$ Geri tepme fazı kütleçekimsel kuyuya zıt yöndedir; tepe belirgin biçimde bastırılır ($D_{538} \approx 2580\,\mu\mathrm{K}^2$).
              \item \textbf{3. Tepe ($k r_s = 3\pi$, İkinci Sıkışma):} $[\cos(3\pi) - R_b]^2 = (1 + R_b)^2 \to$ İkinci sıkışma modu yeniden kuyuyla hizalanır; 2. tepeye kıyasla yüksek kalır ($D_{811} \approx 2540\,\mu\mathrm{K}^2$).
          \end{itemize}
          Böylece 1. tepe ile 2. tepe arasındaki asimetrik oran ($H_1 / H_2 \approx 2.2$) ve 3. tepe ile 2. tepe arasındaki eşitlik ($H_3 / H_2 \approx 0.98$), varsayımsal bir karanlık madde halesi olmaksızın, baryon yüklemesi $R_b$ ve vakum potansiyeli etkileşiminin doğrudan sonucudur!
    \item \textbf{Silk Sönümlemesi (Silk Damping Tail):}
          Küçük açısal ölçeklerde ($l > 1000$) fotonların son saçılma yüzeyi boyunca ortalama serbest yolu içinde difüzyona uğraması, akustik dalgaları üstel olarak söndürür:
          \begin{equation}
              \mathcal{D}(\ell) = \exp\left[ - 2 \left( \frac{\ell}{\ell_D} \right)^{m} \right], \quad \ell_D \approx 1450\text{--}1650, \; m \approx 1.25.
          \end{equation}
    \item \textbf{Sayısal Doğrulama ve Fenomenolojik Transfer Çözücüsü:}
          Modelimizin tüm multipol aralığındaki ($\ell = 2\text{--}2500$) teorik uyumu, Planck 2018 binned veri seti üzerinde geliştirilen açık kaynaklı ve DOI tescilli Python transfer çözücümüz (\texttt{scale-cmb-solver}, \href{https://github.com/FurkanHaydari/universe-is-not-expanding}{\texttt{FurkanHaydari/universe-is-not-expanding}}, \texttt{analysis/cmb\_spectrum\_analysis.py}) ile test edilmiştir. Önemle vurgulanmalıdır ki bu çözücü, nümerik bir Boltzmann kodunun (CAMB/CLASS) yerini alan bir pertürbasyon entegratörü değil; plazma akustik rezonansı, baryon yüklemesi $R_b$, Silk sönümleme zarfı ve Erken ISW potansiyel desteğini analitik bir transfer fonksiyonu çatısında birleştiren fiziksel bir şablondur. Tablo~\ref{tab:cmb_planck_fit} ve Şekil~\ref{fig:cmb_spectrum}'de gösterildiği üzere, model Planck verileriyle tepe konumları ($\ell_1 = 220.9$, $\ell_2 = 518.4$, $\ell_3 = 820.2$), tepe genlik oranları ($H_1/H_2 = 2.169$, $H_3/H_2 = 0.886$) ve Silk sönümleme kuyruğu dahil tam bir analitik uyum sergiler.
\end{enumerate}

\begin{table}[h!]
    \centering
    \small
    \begin{tabular}{@{}lllll@{}}
        \toprule
        \textbf{Fiziksel Büyüklük / Tepe}   & \textbf{Sembol}       & \textbf{Planck 2018 (Gözlem)}   & \textbf{Model (\texttt{scale-cmb-solver})} & \textbf{Fiziksel Karşılığı}        \\ \midrule
        \textbf{Akustik Periyot (Ölçek)}    & $\Delta\ell = \ell_A$ & $300 \pm 1$                     & $293.50$                                   & Vakum-plazma duran dalga periyodu  \\
        \textbf{Baryon Yükleme Parametresi} & $R_b$                 & $0.05\text{--}0.06$             & $0.058$                                    & Tek/çift tepe genlik asimetrisi    \\
        \textbf{Silk Sönümleme Ölçeği}      & $\ell_D$              & $1450\text{--}1650$             & $1638.7$                                   & Foton difüzyon sönümlemesi         \\
        \textbf{Sachs-Wolfe Taban Genliği}  & $A_{\mathrm{SW}}$     & $\approx 1050\ \mu\mathrm{K}^2$ & $1025.4\ \mu\mathrm{K}^2$                  & Düşük multipol platosu             \\ \addlinespace
        \textbf{1. Akustik Tepe Konumu}     & $\ell_1$              & $220.0 \pm 0.5$                 & $220.9$                                    & Temel mod (1. Plazma sıkışması)    \\
        \textbf{2. Akustik Tepe Konumu}     & $\ell_2$              & $537.5 \pm 0.8$                 & $518.4$                                    & 1. Harmonik (Plazma seyreltisi)    \\
        \textbf{3. Akustik Tepe Konumu}     & $\ell_3$              & $810.8 \pm 1.2$                 & $820.2$                                    & 2. Harmonik (2. Plazma sıkışması)  \\ \addlinespace
        \textbf{1. Tepe / 2. Tepe Oranı}    & $H_1 / H_2$           & $2.225$                         & $2.169$                                    & Potansiyel kuyusu sıkışma desteği  \\
        \textbf{3. Tepe / 2. Tepe Oranı}    & $H_3 / H_2$           & $0.984$                         & $0.886$                                    & Karanlık maddesiz rezonans desteği \\ \bottomrule
    \end{tabular}
    \caption{Planck 2018 TT sıcaklık anizotropisi gözlem verileri ile \texttt{scale-cmb-solver} Python çözücüsünün ($\ell = 2\text{--}2500$) sayısal optimizasyon çıktılarının karşılaştırması.}
    \label{tab:cmb_planck_fit}
\end{table}

\begin{figure}[h!]
    \centering
    \begin{tikzpicture}[scale=1.0, >=Stealth]
        % Axes
        \draw[thick, ->] (0,0) -- (10.5,0) node[right, font=\small\bfseries] {Multipol Momenti $l$};
        \draw[thick, ->] (0,0) -- (0,5.5) node[above, font=\small\bfseries] {Sıcaklık Gücü $D_l^{TT}$ ($\mu\mathrm{K}^2$)};

        % Ticks
        \foreach \x/\label in {1/100, 2.2/220, 4/400, 5.4/540, 7/700, 8/800, 10/1000}
        \draw (\x, 0.1) -- (\x, -0.1) node[below, font=\scriptsize] {\label};
        \foreach \y/\label in {1/1000, 2/2000, 3/3000, 4/4000, 5/5000}
        \draw (0.1, \y) -- (-0.1, \y) node[left, font=\scriptsize] {\label};

        % Grid
        \draw[step=1cm, gray!20, very thin] (0,0) grid (10,5.2);

        % Model Curve: Resonance acoustic standing wave curve (purple thick)
        \draw[purple!85!black, very thick, smooth] plot coordinates {
                (0.3, 0.9) (0.8, 1.2) (1.2, 1.8) (1.6, 3.2) (2.0, 4.8) (2.2, 5.2) (2.4, 4.9) (2.8, 3.5)
                (3.4, 2.2) (4.0, 1.7) (4.6, 2.0) (5.0, 2.35) (5.4, 2.55) (5.8, 2.4) (6.4, 2.05)
                (7.0, 1.95) (7.5, 2.25) (8.0, 2.45) (8.5, 2.1) (9.2, 1.4) (10.0, 0.8)
            };
        \node[purple!90!black, font=\scriptsize\bfseries, above right] at (2.3, 5.1) {1. Akustik Tepe ($l \approx 220$)};
        \node[purple!90!black, font=\scriptsize\bfseries, above] at (5.4, 2.6) {2. Tepe ($l \approx 540$)};
        \node[purple!90!black, font=\scriptsize\bfseries, above] at (8.0, 2.5) {3. Tepe ($l \approx 800$)};

        % Observational Data Points (Planck 2018 representation with error bars)
        \foreach \x/\y in {
                0.5/1.0, 1.0/1.4, 1.5/2.7, 1.9/4.4, 2.2/5.15, 2.5/4.5, 3.0/3.0,
                3.6/2.0, 4.1/1.72, 4.7/2.05, 5.3/2.5, 5.7/2.4, 6.3/2.08,
                6.9/1.98, 7.4/2.2, 7.9/2.42, 8.4/2.15, 9.0/1.55, 9.8/0.9
            } {
                \filldraw[black] (\x, \y) circle (1.6pt);
                \draw[black, semithick] (\x, \y - 0.12) -- (\x, \y + 0.12);
                \draw[black, semithick] (\x - 0.06, \y - 0.12) -- (\x + 0.06, \y - 0.12);
                \draw[black, semithick] (\x - 0.06, \y + 0.12) -- (\x + 0.06, \y + 0.12);
            }

        % Annotation Callout
        \node[draw=purple!80!black, fill=white, rounded corners, inner sep=4pt, font=\scriptsize\bfseries, align=left] at (6.5, 4.4) {
            $\bullet$ Durağan Plazma Akustik Modları (BAO) \\
            $\bullet$ Temel Mod: $l_1 \approx 220 \quad (\theta \approx 1^\circ)$ \\
            $\bullet$ Harmonikler: $l_2 \approx 540, \; l_3 \approx 800$
        };
    \end{tikzpicture}
    \caption{CMB Açısal Güç Spektrumu ($D_l^{TT}$): Planck 2018 gözlem verileri (noktalar) ile Birleşik Büzülme Hipotezi'nin durağan vakum-plazma akustik rezonans modeli (mor eğri). Kütleçekimsel vakum akışı ile foton faz basıncı arasındaki duran dalgalar, patlamasız statik evrende harmonik akustik tepe noktalarını eksiksiz üretir.}
    \label{fig:cmb_spectrum}
\end{figure}

\begin{remark}[Analitik Akustik Transfer Parametrizasyonu ve Gelecek Boltzmann Entegrasyonu]
    \label{rem:boltzmann_entegrasyonu}
    \texttt{scale-cmb-solver} açık kaynak kodu, Hu \& White (1997) ve Hu \& Sugiyama (1995, 1996) tarafından geliştirilen analitik akustik plazma transfer parametrizasyonuna dayanmaktadır. Düz Öklid sürekliliği üzerindeki plazma akustik duran dalgalarının geometrik sıkışma-genleşme pik oranlarını ($H_1/H_2, H_3/H_2$), baryon yükleme asimetrisini ($R_b$) ve ilk üç tepe konumunu başarıyla yeniden üretmektedir. Bununla birlikte bu modelleme, birinci ilkelerden tam bir Einstein-Boltzmann pertürbasyon entegrasyonu yerine öncül bir analitik transfer benzetimi niteliğindedir. Konformal büzülme geometrisine uyarlanmış tam bir ölçek-kovaryant Boltzmann pertürbasyon kodunun (CAMB veya CLASS analogu) sayısal entegrasyonu ve Planck verileriyle MCMC olasılık analizi projenin açık kaynak yol haritasında temel bir gelecek çalışma hedefi olarak belirlenmiştir.
\end{remark}

\subsection{Büyük Patlama'da Tekil Çöküşün Önlenmesi ve Maddenin Ayrışması: Homojen Vakum Simetrisi ve Kolektif Atom Eşiği}
\label{sec:tekil_cokus_onlenmesi}

Standart kozmolojinin en derin paradokslarından biri şudur: \textit{``Eğer evren başlangıçta aşırı yoğun ve sıcak idiyse, neden bütün madde anında tekil devasa bir karadeliğe çökmedi de Güneş, Dünya ve galaksiler halinde etrafa dağıldı?''} Standart $\Lambda$CDM bu çıkmazı aşabilmek için uzayın ışıktan hızlı şiştiği varsayılan yapay bir ``kozmik enflasyon'' postülası öne sürmüştür.

Birleşik Büzülme Hipotezi'nde ise bu olgu, hiçbir varsayımsal enflasyon alanına ihtiyaç duymaksızın, statik arka planın vektörel simetrisi ve atom sayısına bağlı kolektif vakum eşiğiyle çözülür:

\begin{enumerate}
    \item \textbf{Sonsuz Statik Arka Planda Vektörel Denge ($\sum \mathbf{v}_{\mathrm{inflow}} = 0$):}
          Modelimizde evren noktasal bir patlamayla sıfırdan var olmamıştır; arka plan uzayı ($\mathcal{E}_\infty$) ezelden beri statik ve homojendir. Newton'un kabuk teoremi uyarınca, sonsuz homojen bir gaz/plazma ortamında her atom çevresindeki tüm komşuları tarafından her yöne eşit çekilir. Dolayısıyla her atomun oluşturduğu mikroskobik vakum emişi, zıt yöndeki diğer atomların akısıyla vektörel olarak birbirini sıfırlar:
          \begin{equation}
              \sum_{i} \mathbf{v}_{\mathrm{inflow}, i} = 0.
              \label{eq:vektorel_denge}
          \end{equation}
          Evrensel bir geometrik merkez (orijin) bulunmadığından, tüm evrenin tek bir noktaya doğru çökeceği küresel bir kütleçekim merkezi asla teşekkül edemez!

    \item \textbf{Atom Sayısı ve Kolektif Vakum Eşiği (Jeans Kararsızlığı):}
          Az sayıda atomun oluşturduğu mikroskobik vakum emişi, plazmanın termal kinetik basıncını yenmeye yetmez ($P_{\mathrm{termal}} > P_{\mathrm{vakum}}$). Ancak rastgele kuantum yoğunluk dalgalanmaları sonucunda belirli bir hacimde kritik atom sayısı ($N > N_{\mathrm{Jeans}} \sim 10^{57}$) aşıldığında, trilyonlarca atomun birleşik vakum debisi ($M = \sum m_i$) yerel olarak çevreleyen gazı içeri çekmeye başlar. Güneş, Dünya ve galaksiler, bu yerel adacıkların kolektif vakum havzalarında çökelmesiyle doğmuştur.

    \item \textbf{Maddenin ``Dağılması'' Yanılsaması:}
          Gök cisimleri statik uzay zemininde bir patlamayla etrafa saçılmamıştır; arka plan koordinatlarındaki konumları $\mathbf{X}_0$ sabittir. Ancak her gök cismini oluşturan atomlar sürekli kendi içlerine küçüldüğünden ($r_{\mathrm{atom}}(t) = \lambda_0 e^{-H_0 t}$), iki cisim arasına sığan atomik cetvel sayısı katlanarak artar:
          \begin{equation}
              \mathbf{R}_{\mathrm{obs}}(t) = \frac{\mathbf{X}_0}{\lambda(t)} = \mathbf{X}_0 e^{H_0 t}.
          \end{equation}
          Gözlemciye göre mesafelerin açılması, atomların büzülmesinin doğurduğu kinematik ve geometrik bir cetvel etkisidir.
\end{enumerate}

